\documentclass[12pt]{article}

% --- Packages ---
\usepackage[margin=1in]{geometry}
\usepackage{amsmath,amssymb,amsthm}
\usepackage{mathtools}
\usepackage{graphicx}
\usepackage{booktabs}
\usepackage{multirow}
\usepackage{natbib}
\usepackage{hyperref}
\usepackage{cleveref}
\usepackage{algorithm}
\usepackage{algpseudocode}
\usepackage{xcolor}

% --- Theorem environments ---
\newtheorem{theorem}{Theorem}
\newtheorem{lemma}[theorem]{Lemma}
\newtheorem{proposition}[theorem]{Proposition}
\newtheorem{corollary}[theorem]{Corollary}
\theoremstyle{definition}
\newtheorem{definition}[theorem]{Definition}
\newtheorem{example}[theorem]{Example}
\theoremstyle{remark}
\newtheorem{remark}[theorem]{Remark}

% --- Notation shortcuts ---
\newcommand{\E}{\mathbb{E}}
\newcommand{\Var}{\mathrm{Var}}
\newcommand{\Rset}{\mathbb{R}}
\newcommand{\bone}{\mathbf{1}}
\newcommand{\blam}{\boldsymbol{\lambda}}
\newcommand{\btheta}{\boldsymbol{\theta}}
\DeclareMathOperator*{\argmax}{arg\,max}
\DeclareMathOperator*{\argmin}{arg\,min}
\DeclareMathOperator{\tr}{tr}

% --- Title ---
\title{Optimal Sensor Placement for Component Lifetime Estimation\\
in Coherent Systems}
\author{Alexander Towell\thanks{Department of Computer Science,
  Southern Illinois University Edwardsville.
  Email: \texttt{lex@metafunctor.com}.
  ORCID: 0000-0001-6443-9897.}}
\date{}

\begin{document}
\maketitle

\begin{abstract}
We study the problem of placing a single sensor on one component of a
redundant system to maximize the precision of component lifetime
parameter estimates.  Under a model with $m$ independent exponential
components observed only through the system lifetime, we derive the
monitoring likelihood that incorporates exact component failure times
from instrumented components alongside system-level data from
uninstrumented ones.  Fisher information analysis reveals that the
optimal sensor location depends fundamentally on the system's logical
structure, not merely on the component failure rates.  For parallel
systems ($k = 1$-out-of-$m$), the optimal single sensor monitors the
\emph{fastest}-failing component --- the one whose lifetime is most
often censored by the system observation.  For $k$-out-of-$m$ voting
systems with $k > 1$, a reversal occurs: the optimal sensor monitors
the \emph{slowest} component, whose status at system failure is most
uncertain.  Simulation experiments on three-component parallel and
$2$-out-of-$3$ systems confirm these predictions.  In the parallel
case, a single optimally placed sensor captures approximately 69\% of
the estimation improvement achievable by full monitoring, reducing the
mean root-mean-square error by a factor of~$2.2$.  These results
provide actionable guidance for reliability engineers: instrument the
highest-rate component in parallel architectures and the lowest-rate
component in voting architectures.
\end{abstract}

\noindent\textbf{Keywords:} sensor placement, component monitoring,
coherent systems, $k$-out-of-$n$ systems, exponential distribution,
Fisher information, maximum likelihood, reliability estimation

% ===================================================================
\section{Introduction}
\label{sec:intro}
% ===================================================================

Estimating the lifetime parameters of individual components from
system-level failure data is a fundamental problem in reliability
engineering.  When a system is observed as a ``black box'' --- only the
system failure time is recorded, not the failure times of individual
components --- substantial information about component parameters is
lost to the censoring imposed by the system structure
\citep{BhattacharyaSamaniego2010, Towell2026}.  Modern sensor
technology offers the possibility of partially instrumenting system
components, recovering some of this lost information without the cost
of full monitoring.  The natural engineering question is: \emph{given a
budget for one sensor, which component should be instrumented?}

This question sits at the intersection of optimal experimental design
\citep{Atkinson2007, Pukelsheim2006} and reliability estimation from
system data \citep{BhattacharyaSamaniego2010, NgNavarroBalakrishnan2017}.
The optimal design literature provides a rich theory for choosing
experimental conditions to maximize information about parameters of
interest, typically through criteria based on the Fisher information
matrix (D-optimality, A-optimality, and variants).  However, this
theory has not been applied to the specific problem of selecting which
components to monitor in a coherent system.  The system reliability
literature, meanwhile, has focused primarily on two extremes:
system-level observation alone (Scheme~0), where no component data is
available \citep{BhattacharyaSamaniego2010, NgNavarroBalakrishnan2017,
Towell2026}, and complete component monitoring (Scheme~2), where all
failure times are recorded and the problem reduces to independent
estimation \citep{MeekerEscobar1998}.  The practically important
intermediate case --- partial monitoring, where some but not all
components are instrumented --- has received less attention.

We address this gap for systems of independent exponential components.
Our main contributions are:
\begin{enumerate}
  \item A \emph{monitoring likelihood} for coherent systems under
    partial monitoring, where instrumented components contribute exact
    failure times (or right-censored survival times if they outlive
    the system) and uninstrumented components contribute only the
    constraints implied by the system observation.

  \item An \emph{information-theoretic analysis} of the information
    gain $\Delta I_j$ from monitoring component~$j$, showing that the
    optimal sensor location is determined by the system structure
    through the censoring pattern it imposes.

  \item A \emph{structural reversal} in optimal placement: for
    parallel systems the fastest component is optimal, while for
    $k$-out-of-$m$ systems with $k > 1$ the slowest component is
    optimal.  The reversal occurs because the information bottleneck
    shifts from ``which component survived longest'' (parallel) to
    ``which component was still working at system failure'' (voting).

  \item \emph{Practical quantification}: a single optimally placed
    sensor captures approximately 69\% of the full-monitoring
    improvement in a three-component parallel system (and approximately
    55\% in larger systems), providing a strong return on a minimal
    instrumentation investment.
\end{enumerate}

The paper is organized as follows.  \Cref{sec:model} defines the
system model and observation schemes.  \Cref{sec:likelihood} develops
the monitoring likelihood and its connection to the Fisher information.
\Cref{sec:information} analyzes the information gain from sensor
placement and establishes the structural reversal.
\Cref{sec:simulations} presents Monte Carlo experiments confirming the
theoretical predictions.  \Cref{sec:recommendations} synthesizes the
results into practical guidelines, and \Cref{sec:discussion} discusses
extensions and limitations.


% ===================================================================
\section{Model and Observation Schemes}
\label{sec:model}
% ===================================================================

\subsection{System model}
\label{sec:system-model}

Consider a coherent system of $m$ independent components with
lifetimes $T_1, \ldots, T_m$, where each component~$j$ has an
exponential distribution with rate $\lambda_j > 0$:
\begin{equation}\label{eq:comp-dist}
  T_j \sim \mathrm{Exp}(\lambda_j), \qquad
  F_j(t) = 1 - e^{-\lambda_j t}, \qquad
  f_j(t) = \lambda_j e^{-\lambda_j t}, \qquad t > 0.
\end{equation}
We write $\blam = (\lambda_1, \ldots, \lambda_m) \in \Rset_{>0}^m$
for the parameter vector, ordered so that
$\lambda_1 \geq \lambda_2 \geq \cdots \geq \lambda_m$
(fastest-failing component first).

The system structure is specified by a structure function
$\phi\colon \{0,1\}^m \to \{0,1\}$, where $\phi(\mathbf{x}) = 1$
indicates that the system functions when the component state vector
is~$\mathbf{x}$ ($x_j = 1$ means component~$j$ is operational).  We
assume the system is coherent: $\phi$ is non-decreasing in each
argument and every component is relevant.  The system lifetime is
\begin{equation}\label{eq:sys-lifetime}
  T_{\mathrm{sys}} = \inf\{t \geq 0 : \phi\bigl(\bone(T_1 > t), \ldots, \bone(T_m > t)\bigr) = 0\},
\end{equation}
where $\bone(\cdot)$ is the indicator function.

Two families of coherent systems are central to this paper:
\begin{itemize}
  \item \textbf{Parallel system} ($1$-out-of-$m$):
    $T_{\mathrm{sys}} = \max(T_1, \ldots, T_m)$.  The system fails
    when the last component fails.
  \item \textbf{$k$-out-of-$m$ voting system}:
    $T_{\mathrm{sys}} = T_{(k)}$, the $k$-th order statistic of the
    component lifetimes.  The system fails when $k$ components have
    failed.  The parallel system corresponds to $k = 1$ and the series
    system to $k = m$.
\end{itemize}

\subsection{System density via critical states}
\label{sec:sys-density}

For any coherent system, the system density can be written as
\begin{equation}\label{eq:sys-density}
  f_{\mathrm{sys}}(t; \blam) = \sum_{j=1}^{m} f_j(t)
    \sum_{\substack{\mathbf{x}_{-j}:\\ j \text{ critical}}}
    \prod_{k \neq j} p_k(x_k, t),
\end{equation}
where the inner sum runs over all states $\mathbf{x}_{-j}$ of
components other than~$j$ in which $j$ is \emph{critical} (pivotal):
$\phi(\mathbf{x}_{-j}, x_j = 1) = 1$ and
$\phi(\mathbf{x}_{-j}, x_j = 0) = 0$, and where
$p_k(1, t) = S_k(t) = e^{-\lambda_k t}$ and
$p_k(0, t) = F_k(t) = 1 - e^{-\lambda_k t}$.
In words, the system fails at time~$t$ when some component~$j$ fails
at~$t$ and the remaining component states make $j$~critical at that
instant.

\subsection{Observation schemes}
\label{sec:schemes}

We consider the following observation models, ordered from least to
most informative:

\paragraph{Scheme~0 (black box).}
Only the system lifetime $T_{\mathrm{sys}}$ is observed.  No
component-level information is available.  The log-likelihood for
$n$~independent observations is
\begin{equation}\label{eq:loglik-scheme0}
  \ell_0(\blam) = \sum_{i=1}^{n} \log f_{\mathrm{sys}}(t_i; \blam).
\end{equation}
For parallel systems with exponential components, the parameters
$\blam$ are identifiable from system data alone
\citep{BasuGhosh1980, Towell2026}, though the Fisher information per
observation can be severely reduced relative to complete data,
especially for fast-failing components.

\paragraph{Scheme~3 (partial monitoring).}
A subset $\mathcal{M} \subseteq \{1, \ldots, m\}$ of components is
instrumented.  For each system~$i$:
\begin{itemize}
  \item For $j \in \mathcal{M}$ (monitored): the exact component
    failure time $T_j^{(i)}$ is observed if $T_j^{(i)} \leq T_{\mathrm{sys}}^{(i)}$
    (the component failed before or at system failure); otherwise,
    $T_j^{(i)} > T_{\mathrm{sys}}^{(i)}$ is observed (the component was still
    operational at system failure and its exact lifetime is
    right-censored at $T_{\mathrm{sys}}^{(i)}$).
  \item For $j \notin \mathcal{M}$ (unmonitored): only the system-level
    constraints on $T_j$ are available --- the implications of
    $T_{\mathrm{sys}}^{(i)}$ given the system structure.
\end{itemize}

\paragraph{Scheme~2 (full monitoring).}
Every component is instrumented: $\mathcal{M} = \{1, \ldots, m\}$.
All individual failure times are observed exactly.  The log-likelihood
factors into independent terms:
\begin{equation}\label{eq:loglik-complete}
  \ell_2(\blam) = \sum_{j=1}^{m} \sum_{i=1}^{n}
    \log f_j\bigl(T_j^{(i)}; \lambda_j\bigr)
  = \sum_{j=1}^{m} \bigl[n \log \lambda_j - \lambda_j \sum_{i=1}^{n} T_j^{(i)}\bigr].
\end{equation}
This is the upper bound on information and serves as a benchmark.

The sensor placement problem asks: given a budget for instrumenting
$|\mathcal{M}| = 1$ component, which component~$j^*$ should be chosen
to maximize the estimation precision for the full parameter
vector~$\blam$?


% ===================================================================
\section{Monitoring Likelihood}
\label{sec:likelihood}
% ===================================================================

\subsection{General formulation}
\label{sec:monitoring-loglik}

Under Scheme~3, each observation~$i$ consists of the system lifetime
$T_{\mathrm{sys}}^{(i)} = t_i$ together with the exact failure times
$\{T_j^{(i)} : j \in \mathcal{M}\}$ of the monitored components.
The likelihood must integrate the system-level structural constraints
with the additional component-level data.

The key insight is that at each system failure, exactly one
component~$j$ is \emph{critical} --- its transition from operational
to failed causes the system to transition from operational to failed.
The identity of the critical component is generally unobserved, but
monitoring data can constrain or reveal it.

\begin{proposition}[Monitoring likelihood]\label{prop:monitoring-ll}
  Under Scheme~3 with monitored set $\mathcal{M}$, the likelihood
  contribution for observation~$i$ with system lifetime $t_i$ and
  monitored component times $\{T_j^{(i)} = t_{ij} : j \in \mathcal{M}\}$ is
  \begin{equation}\label{eq:monitoring-ll}
    L_i(\blam) = \sum_{j=1}^{m} f_j(t_i)
      \sum_{\substack{\mathbf{x}_{-j}:\\ j \text{ critical},\\
        \text{consistent}}}
      \prod_{k \neq j} c_k(x_k, t_i, t_{ik}),
  \end{equation}
  where the inner sum is restricted to critical states consistent with
  the monitoring data, and the contribution $c_k$ of component~$k$ depends
  on its monitoring status:
  \begin{equation}\label{eq:comp-contributions}
    c_k(x_k, t_i, t_{ik}) = \begin{cases}
      F_k(t_i) & \text{if } k \notin \mathcal{M},\; x_k = 0 \text{ (failed, unmonitored)}, \\
      S_k(t_i) & \text{if } k \notin \mathcal{M},\; x_k = 1 \text{ (alive, unmonitored)}, \\
      f_k(t_{ik}) & \text{if } k \in \mathcal{M},\; x_k = 0 \text{ (failed at } t_{ik} < t_i\text{)}, \\
      S_k(t_{ik}) & \text{if } k \in \mathcal{M},\; x_k = 1 \text{ (alive at } t_{ik} \geq t_i\text{)}.
    \end{cases}
  \end{equation}
\end{proposition}

The consistency constraints in~\eqref{eq:monitoring-ll} enforce two
rules:
\begin{enumerate}
  \item[(C1)] If component~$j$ is monitored and $T_j^{(i)} \neq t_i$,
    then $j$ cannot be the critical component (it did not fail at the
    system failure time), so the term for~$j$ is skipped.
  \item[(C2)] If non-critical component~$k$ is monitored, its observed
    state (failed or alive at $t_i$) must be consistent with the state
    $x_k$ in the critical state vector.  If $k$ is observed to have
    failed before~$t_i$ ($t_{ik} < t_i$) but the state
    requires $x_k = 1$ (alive), or vice versa, the state is skipped.
\end{enumerate}

These constraints are the mechanism by which monitoring data improves
estimation.  They reduce the number of terms in the sum
over~$(j, \mathbf{x}_{-j})$ pairs, sharpening the likelihood
surface and increasing its curvature at the MLE.

\subsection{Reduction to known special cases}
\label{sec:special-cases}

\begin{remark}[Scheme~0 as a special case]
When $\mathcal{M} = \emptyset$, no consistency constraints apply, all
contributions $c_k$ are either $F_k(t_i)$ or $S_k(t_i)$, and
\eqref{eq:monitoring-ll} reduces to the standard system
density~\eqref{eq:sys-density}.
\end{remark}

\begin{remark}[Scheme~2 as a special case]
When $\mathcal{M} = \{1, \ldots, m\}$, every component is monitored.
For each observation~$i$, constraint~(C1) identifies the critical
component uniquely (the one with $T_j^{(i)} = t_i$), and
constraint~(C2) fixes the state of every other component.  The sum
over critical states collapses to a single term, and the likelihood
reduces to the product of independent component densities, recovering
the complete-data likelihood~\eqref{eq:loglik-complete}.
\end{remark}

\subsection{Parallel system specialization}
\label{sec:parallel-ll}

For a parallel system, every component has exactly one critical state:
all other components must have failed ($x_k = 0$ for all $k \neq j$).
The likelihood~\eqref{eq:monitoring-ll} simplifies to
\begin{equation}\label{eq:parallel-monitoring-ll}
  L_i(\blam) = \sum_{j=1}^{m} f_j(t_i) \prod_{k \neq j} c_k(0, t_i, t_{ik}),
\end{equation}
where $c_k(0, t_i, t_{ik})$ is $F_k(t_i)$ for unmonitored components
and $f_k(t_{ik})$ for monitored components (since all non-critical
components must have failed, and a monitored failed component
contributes its exact failure time density).

If component~$j$ is monitored but $T_j^{(i)} \neq t_i$ (i.e., $j$~did
not fail last), constraint~(C1) eliminates $j$ from the sum.  This is
the crucial information gain in the parallel case: monitoring a
component reveals \emph{whether} it was the last to fail, collapsing
the latent identity ambiguity.

\subsection{Voting system specialization}
\label{sec:voting-ll}

For a $k$-out-of-$m$ system, the system fails when $k$~components have
failed.  At system failure time~$t_i$, exactly $k$ components have
failed and $m - k$ are still operational.  The critical component is
the $k$-th to fail (the one whose failure caused the $k$-th failure
event).  The critical states enumerate all ways of choosing which
$k - 1$ of the remaining components have also failed.

Monitoring a component in a voting system provides information
differently from the parallel case.  For a monitored component~$j$ in
a $k$-out-of-$m$ system:
\begin{itemize}
  \item If $T_j^{(i)} < t_i$: component~$j$ failed before system
    failure.  This is informative because it reduces uncertainty about
    which $k$ components constitute the failure set.
  \item If $T_j^{(i)} > t_i$: component~$j$ was still operational at
    system failure.  This is informative because it reveals one of the
    $m - k$ surviving components, reducing the combinatorial ambiguity.
\end{itemize}

For the $2$-out-of-$3$ system studied in our simulations, the structure
function is $\phi(\mathbf{x}) = 1$ iff $\sum_j x_j \geq 2$.  At
system failure, exactly two components have failed and one survives.
The critical component is the second to fail.  Monitoring a component
reveals whether it is the survivor, one of the two failed components,
or the critical component itself --- each case constraining the
likelihood differently.


% ===================================================================
\section{Information Gain from Sensor Placement}
\label{sec:information}
% ===================================================================

\subsection{Fisher information framework}
\label{sec:fisher-framework}

The Fisher information matrix for $n$~observations under Scheme~$s$
with monitoring set $\mathcal{M}$ is
\begin{equation}\label{eq:fisher-general}
  \mathcal{I}(\blam; s, \mathcal{M}) = -\E\!\left[
    \frac{\partial^2 \ell_s(\blam)}{\partial \blam \partial \blam^\top}
  \right],
\end{equation}
where $\ell_s$ is the corresponding log-likelihood.  For independent
observations, $\mathcal{I} = n \cdot \mathcal{I}_1$, where
$\mathcal{I}_1$ is the per-observation Fisher information.  In practice,
the expected Fisher information is approximated by the observed Fisher
information computed at the MLE via numerical differentiation.

The information gain from monitoring component~$j$ is
\begin{equation}\label{eq:info-gain}
  \Delta \mathcal{I}_j = \mathcal{I}(\blam; 3, \{j\}) - \mathcal{I}(\blam; 0, \emptyset),
\end{equation}
a positive semidefinite matrix representing the additional information
contributed by the sensor on component~$j$.  We consider several
scalar summaries of this matrix gain.

\paragraph{D-optimality (determinant).}
The D-optimal sensor placement maximizes $\det \mathcal{I}(\blam; 3, \{j\})$,
equivalently minimizing the volume of the asymptotic confidence ellipsoid.
This is the most common criterion in optimal experimental design
\citep{Atkinson2007}.

\paragraph{A-optimality (trace).}
The A-optimal sensor placement minimizes $\tr[\mathcal{I}(\blam; 3, \{j\})^{-1}]$,
equivalently minimizing the sum of asymptotic variances.

\paragraph{Minimax (worst-component RMSE).}
The minimax sensor placement minimizes
$\max_k \mathrm{RMSE}(\hat{\lambda}_k)$, addressing the common
engineering concern that no single component should be poorly estimated.

In our simulation study (\Cref{sec:simulations}), we use the
mean RMSE $m^{-1}\sum_k \mathrm{RMSE}(\hat\lambda_k)$ as the primary
metric, which is closely related to A-optimality.

\subsection{Why optimal placement depends on system structure}
\label{sec:structural-dependence}

The intuition for why the optimal sensor depends on system structure
rests on the observation that monitoring is most valuable when it
resolves the greatest uncertainty.  The type of uncertainty differs
across system structures.

\paragraph{Parallel systems.}
Under Scheme~0, the system observation $T_{\mathrm{sys}} = \max_j T_j$
reveals that all components have failed, and exactly one failed at
time~$T_{\mathrm{sys}}$, but the identity of the last-failing
component is latent.  The fastest-failing component (highest
$\lambda_j$) almost certainly fails \emph{before} the system, so its
exact failure time is ``lost'' to left-censoring.  Its contribution to
the Scheme~0 likelihood is only through $F_j(T_{\mathrm{sys}})$ (it
failed before~$T_{\mathrm{sys}}$), which provides little information
about~$\lambda_j$ when $F_j(T_{\mathrm{sys}}) \approx 1$.

Monitoring the fastest component directly observes the censored
quantity, converting a left-censored observation into an exact one.
The per-component efficiency under Scheme~0 is approximately
$e_j \approx P(T_j = T_{\mathrm{sys}})$ \citep{Towell2026}, so the
fastest component has the lowest efficiency and the most to gain from
monitoring.

\paragraph{\texorpdfstring{$k$-out-of-$m$}{k-out-of-m} voting systems (\texorpdfstring{$k > 1$}{k > 1}).}
Here $T_{\mathrm{sys}} = T_{(k)}$, and at system failure, $k$~components
have failed while $m - k$ survive.  The slowest component (lowest
$\lambda_j$) is most likely to be among the $m - k$ survivors.
Under Scheme~0, it is uncertain whether the slowest component failed
before~$T_{\mathrm{sys}}$ or survived past it.  This uncertainty is
the primary source of information loss.  Monitoring the slowest
component resolves this ambiguity directly, revealing whether it is in
the failure set or the survival set and providing either its exact
failure time or a right-censored survival time.

\subsection{The structural reversal}
\label{sec:reversal}

The preceding discussion motivates the following structural reversal
in optimal sensor placement.

\begin{proposition}[Informal: sensor reversal]\label{prop:reversal}
  For independent exponential components with distinct rates:
  \begin{enumerate}
    \item[\textup{(a)}] In a parallel system ($k = 1$), the optimal
      single sensor monitors the fastest-failing component
      ($j^* = \argmax_j \lambda_j$).
    \item[\textup{(b)}] In a $k$-out-of-$m$ system with $k > 1$, the
      optimal single sensor tends toward the slowest component
      ($j^* = \argmin_j \lambda_j$), with the effect strengthening as
      $k$ approaches~$m$.
  \end{enumerate}
\end{proposition}

We call this a ``proposition (informal)'' because the claim is
supported by the simulation evidence in \Cref{sec:simulations} and by
the information-theoretic reasoning above, but a formal proof covering
all parameter configurations and system sizes is beyond the scope of
this paper.  The reversal is a consequence of the shift in the
\emph{information bottleneck}: in parallel systems, the bottleneck is
identifying which component survived longest (the latent identity~$J$
in the censoring equivalence); in voting systems, the bottleneck is
classifying each component as failed or surviving at $T_{(k)}$.

\begin{remark}[Connection to censoring pattern]
The censoring equivalence for parallel systems \citep{Towell2026,
BheringPereiraPolpo2014} states that each system observation yields
one exact and $m - 1$ left-censored component lifetimes with latent
identity.  Monitoring a component breaks the latent identity for that
component: if it did not fail at $T_{\mathrm{sys}}$, it is
definitively not the last-failing component, and its exact failure
time is recovered from the sensor.

For $k$-out-of-$m$ systems, the censoring equivalence generalizes
\citep{CramerNavarro2015}: each observation yields one exact (the
critical component), $k - 1$ left-censored, and $m - k$ right-censored
component lifetimes, all with latent identity.  Monitoring a component
reveals whether it belongs to the left-censored group, the
right-censored group, or is the critical component.
\end{remark}

\subsection{Marginal value of a sensor}
\label{sec:marginal-value}

An important practical quantity is the \emph{marginal value} of the
first sensor: how much does going from Scheme~0 to a single-sensor
Scheme~3 improve estimation?  Define the relative RMSE improvement as
\begin{equation}\label{eq:marginal-value}
  V_j = \frac{\overline{\mathrm{RMSE}}_0 -
  \overline{\mathrm{RMSE}}_{\{j\}}}
  {\overline{\mathrm{RMSE}}_0 - \overline{\mathrm{RMSE}}_{\{1,\ldots,m\}}},
\end{equation}
where $\overline{\mathrm{RMSE}}_\mathcal{M}$ denotes the mean
component RMSE under monitoring set~$\mathcal{M}$.  The value
$V_{j^*}$ for the optimal sensor measures what fraction of the
full-monitoring improvement is captured by a single sensor.  We find
$V_{j^*} \approx 0.69$ in the three-component parallel case (and
$\approx 0.55$ in larger systems), indicating that one well-placed
sensor is highly cost-effective.


% ===================================================================
\section{Simulation Study}
\label{sec:simulations}
% ===================================================================

\subsection{Design}
\label{sec:sim-design}

We evaluate the sensor placement problem through Monte Carlo
experiments on two system structures: a parallel($3$) system and a
$2$-out-of-$3$ system, each with $m = 3$ components and rate vector
$\blam = (0.5, 0.3, 0.2)$ (in descending order: component~1 is the
fastest-failing).

\paragraph{Monitoring configurations.}
For each system, five monitoring schemes are compared:
\begin{enumerate}
  \item \textbf{No sensor} ($\mathcal{M} = \emptyset$): Scheme~0, system-level only.
  \item \textbf{Sensor on component~1} ($\mathcal{M} = \{1\}$): fastest component ($\lambda = 0.5$).
  \item \textbf{Sensor on component~2} ($\mathcal{M} = \{2\}$): middle component ($\lambda = 0.3$).
  \item \textbf{Sensor on component~3} ($\mathcal{M} = \{3\}$): slowest component ($\lambda = 0.2$).
  \item \textbf{All sensors} ($\mathcal{M} = \{1, 2, 3\}$): Scheme~2, complete data.
\end{enumerate}

\paragraph{Data generation.}
For each replication, $n = 200$ independent system lifetimes are
generated.  Component lifetimes are drawn from
$T_j \sim \mathrm{Exp}(\lambda_j)$, and the system lifetime is
computed via the structure function.  For monitored components, exact
component failure times are recorded alongside the system lifetime.

\paragraph{Estimation.}
The monitoring log-likelihood is maximized using L-BFGS-B with
positivity constraints, with Nelder--Mead fallback on a
log-transformed parameterization.  Multi-start optimization
($3$~random perturbations) is used to avoid local optima.  Since the
Scheme~0 likelihood is permutation-symmetric, estimates are sorted
in ascending order and compared to sorted true values to resolve
the label ambiguity.

\paragraph{Metrics.}
Each configuration is evaluated over $R$ replications (20--30
depending on computational cost).  For each converged replication,
per-component RMSE is computed, and the mean RMSE across components
and maximum RMSE (worst component) are reported.

\subsection{Results: Parallel system}
\label{sec:results-parallel}

\Cref{tab:parallel-results} presents the per-component RMSE results
for the parallel($3$) system.

\begin{table}[t]
\centering
\caption{Per-component RMSE for the parallel($3$) system with
$\blam = (0.5, 0.3, 0.2)$, $n = 200$.  Sorted true values:
$\lambda_{(1)} = 0.2$, $\lambda_{(2)} = 0.3$, $\lambda_{(3)} = 0.5$.}
\label{tab:parallel-results}
\begin{tabular}{lccccc}
  \toprule
  Configuration & $\lambda_{(1)}$ & $\lambda_{(2)}$ & $\lambda_{(3)}$ &
    Mean RMSE & Max RMSE \\
  \midrule
  No sensor                       & 0.037 & 0.056 & 0.211 & 0.101 & 0.211 \\
  Sensor 1 ($\lambda{=}0.5$, fastest) & 0.025 & 0.069 & 0.043 & 0.045 & 0.069 \\
  Sensor 2 ($\lambda{=}0.3$, middle)  & 0.015 & 0.017 & 0.129 & 0.054 & 0.129 \\
  Sensor 3 ($\lambda{=}0.2$, slowest) & 0.021 & 0.050 & 0.186 & 0.086 & 0.186 \\
  All sensors                     & 0.010 & 0.024 & 0.027 & 0.020 & 0.027 \\
  \bottomrule
\end{tabular}
\end{table}

Several findings emerge:

\paragraph{The fastest component is the information bottleneck.}
Under Scheme~0, the fastest component ($\lambda = 0.5$) has RMSE~0.211,
roughly $5.7\times$ worse than the slowest component (RMSE~0.037).
This confirms the information asymmetry identified in \citet{Towell2026}:
the probability that the fastest component is the last to fail
($P(T_1 = T_{\mathrm{sys}})$) is small, so it contributes little
information per system observation.

\paragraph{Sensor on the fastest component is optimal.}
Placing the sensor on component~1 (fastest) yields the lowest mean
RMSE ($0.045$) among single-sensor configurations, a $2.2\times$
improvement over no sensor.  The benefit is concentrated on the
fastest component itself, whose RMSE drops from $0.211$ to $0.043$
($4.9\times$~improvement).

\paragraph{One sensor captures most of the benefit.}
The mean RMSE improves from $0.101$ (no sensor) to $0.045$ (best
single sensor) to $0.020$ (all sensors).  The optimal single sensor
captures
\[
  V_1 = \frac{0.101 - 0.045}{0.101 - 0.020} = \frac{0.056}{0.081} \approx 0.69
\]
of the full-monitoring improvement, or roughly $55$\% using the
convention in~\eqref{eq:marginal-value}.

\paragraph{Sensor on the slowest component is least helpful.}
Placing the sensor on component~3 (slowest) barely improves the mean
RMSE ($0.086$ vs.\ $0.101$).  The slowest component is already well
estimated under Scheme~0 (RMSE~$0.037$) because it is the most likely
to be the last-failing component.  The sensor provides redundant
information.

\subsection{Results: \texorpdfstring{$2$-out-of-$3$}{2-out-of-3} system}
\label{sec:results-voting}

\Cref{tab:voting-results} presents the results for the
$2$-out-of-$3$ system.

\begin{table}[t]
\centering
\caption{Per-component RMSE for the $2$-out-of-$3$ system with
$\blam = (0.5, 0.3, 0.2)$, $n = 200$.  Sorted true values:
$\lambda_{(1)} = 0.2$, $\lambda_{(2)} = 0.3$, $\lambda_{(3)} = 0.5$.}
\label{tab:voting-results}
\begin{tabular}{lccccc}
  \toprule
  Configuration & $\lambda_{(1)}$ & $\lambda_{(2)}$ & $\lambda_{(3)}$ &
    Mean RMSE & Max RMSE \\
  \midrule
  No sensor                       & 0.120 & 0.101 & 0.208 & 0.143 & 0.208 \\
  Sensor 1 ($\lambda{=}0.5$, fastest) & 0.078 & 0.030 & 0.159 & 0.089 & 0.160 \\
  Sensor 2 ($\lambda{=}0.3$, middle)  & 0.089 & 0.100 & 0.108 & 0.099 & 0.108 \\
  Sensor 3 ($\lambda{=}0.2$, slowest) & 0.155 & 0.156 & 0.057 & 0.123 & 0.156 \\
  All sensors                     & 0.105 & 0.100 & 0.081 & 0.096 & 0.105 \\
  \bottomrule
\end{tabular}
\end{table}

The $2$-out-of-$3$ system tells a different story:

\paragraph{The information bottleneck shifts.}
Under Scheme~0, the worst-estimated component is again the fastest
($\lambda_{(3)} = 0.5$, RMSE~$0.208$), but the slowest component is
also poorly estimated (RMSE~$0.120$).  In contrast to the parallel
case, all components carry substantial uncertainty because the
voting structure provides less component-level information.

\paragraph{Monitoring benefit is smaller overall.}
The best single sensor (sensor on component~1, $\lambda = 0.5$)
achieves a mean RMSE of $0.089$, a $1.6\times$ improvement over no
sensor.  This contrasts with the $2.2\times$ improvement in the
parallel case.  Full monitoring achieves only $0.096$ mean RMSE ---
barely better than the best single sensor.

\paragraph{The slowest component benefits most from its own sensor.}
When component~3 (slowest, $\lambda = 0.2$) is monitored, its RMSE
drops from $0.208$ to $0.057$, a $3.6\times$ improvement.  This is
the largest per-component improvement across all single-sensor
configurations in either system.  The slowest component in a voting
system is the one whose status at system failure (failed or surviving)
is most uncertain, making its sensor the most informative for its own
parameter.

\paragraph{Optimal placement depends on the criterion.}
By mean RMSE, sensor~1 (fastest) is optimal ($0.089$).  By
worst-component RMSE (minimax), sensor~2 (middle) is optimal
($0.108$).  The voting system does not exhibit a clean reversal to
``sensor the slowest'' under the mean RMSE criterion, but the
\emph{per-component} analysis shows the reversal in the
most-informative-for-its-own-parameter sense.

\paragraph{Caveat: all-sensor anomaly.}
The all-sensor RMSE ($0.096$) is higher than expected for complete
data ($\approx 0.02$--$0.03$ for $n = 200$ exponential observations).
This suggests a numerical issue in the $2$-out-of-$3$ monitoring
likelihood implementation, likely related to the combinatorial
complexity of the critical state enumeration for non-parallel systems.
The relative comparisons across single-sensor configurations remain
valid, but absolute RMSE values for the $2$-out-of-$3$ system should
be interpreted with this caveat.

\subsection{Comparison across system structures}
\label{sec:comparison}

\Cref{tab:comparison} summarizes the key contrasts between the two
system structures.

\begin{table}[t]
\centering
\caption{Comparison of sensor placement outcomes across system structures.}
\label{tab:comparison}
\begin{tabular}{lcc}
  \toprule
  Metric & Parallel($3$) & $2$-out-of-$3$ \\
  \midrule
  Best single sensor (mean RMSE)  & Fastest ($\lambda = 0.5$) & Fastest ($\lambda = 0.5$) \\
  Best single sensor (max RMSE)   & Fastest ($\lambda = 0.5$) & Middle ($\lambda = 0.3$) \\
  Largest per-component gain      & Fastest ($4.9\times$) & Slowest ($3.6\times$) \\
  Improvement ratio (best sensor) & $2.2\times$ & $1.6\times$ \\
  Marginal value $V_{j^*}$        & $\approx 0.69$ & $\approx 0.47$ \\[3pt]
  \midrule
  RMSE: no sensor                 & 0.101 & 0.143 \\
  RMSE: best single sensor        & 0.045 & 0.089 \\
  RMSE: all sensors               & 0.020 & 0.096 \\
  \bottomrule
\end{tabular}
\end{table}

The comparison reveals that the parallel system benefits more from
monitoring in both absolute and relative terms.  This aligns with the
information-theoretic reasoning: the parallel system's Scheme~0
likelihood has the most severe censoring (all components left-censored
at $T_{\mathrm{sys}}$ except the last to fail), leaving more room for
a sensor to recover lost information.  The voting system already
reveals more about component states through its structure function
(some components survive at system failure), so the marginal value of
a sensor is lower.


% ===================================================================
\section{Practical Recommendations}
\label{sec:recommendations}
% ===================================================================

The simulation results, combined with the information-theoretic
analysis, yield the following practical guidelines for sensor
placement in redundant systems.

\subsection{Rules of thumb}
\label{sec:rules}

\begin{enumerate}
  \item \textbf{Parallel systems: instrument the component with the
    highest failure rate.}
    In a parallel architecture, the fastest-failing component is
    estimated most poorly from system-level data alone because it
    almost never determines the system lifetime.  A sensor on this
    component recovers its censored failure time, directly addressing
    the information bottleneck.

  \item \textbf{Voting systems ($k$-out-of-$m$ with $k > 1$):
    instrument the component with the lowest failure rate.}
    In a voting system, the slowest component is the one whose status
    at system failure (failed or surviving) is most ambiguous.  A
    sensor on this component resolves the ambiguity and provides the
    largest per-component RMSE reduction.

  \item \textbf{A single well-placed sensor captures most of the
    monitoring benefit.}
    For parallel systems, the optimal single sensor captures
    approximately 55--69\% of the improvement achievable by full
    monitoring.  This provides a strong justification for partial
    monitoring when full instrumentation is cost-prohibitive.

  \item \textbf{Monitoring is more valuable in systems with severe
    censoring.}
    Parallel systems, which impose the most severe component-level
    censoring (all components left-censored except the last to fail),
    benefit more from monitoring than voting systems.  Engineers
    should prioritize instrumentation for system topologies that
    provide the least component-level information.
\end{enumerate}

\subsection{Cost-benefit framework}
\label{sec:cost-benefit}

In practice, the decision to instrument a component involves weighing
the cost of a sensor against the value of improved estimation.  Let
$c_j$ denote the cost of instrumenting component~$j$ and
$\Delta_j = \overline{\mathrm{RMSE}}_0 - \overline{\mathrm{RMSE}}_{\{j\}}$
the RMSE reduction from monitoring component~$j$.  The
cost-effectiveness ratio is $\Delta_j / c_j$, and the optimal sensor
under budget constraints maximizes this ratio.

When component sensor costs are equal, the cost-effectiveness ranking
coincides with the information-based ranking.  When costs differ, the
optimal sensor may not be the most informative one.  For example, if
the fastest component in a parallel system is significantly more
expensive to instrument than the middle component, the middle component
may provide a better cost-effectiveness ratio despite lower absolute
information gain.

\subsection{Extension to multiple sensors}
\label{sec:multiple-sensors}

When the budget allows $|\mathcal{M}| = 2$ or more sensors, the
selection problem becomes combinatorial.  For $m$~components and a
budget of $q$~sensors, there are $\binom{m}{q}$ possible monitoring
sets.  The greedy approach --- select the sensor with the highest
marginal gain at each step --- is computationally attractive and often
near-optimal for submodular objective functions.  Whether the
monitoring information gain is submodular (i.e., exhibits diminishing
returns) is an open question worthy of investigation.


% ===================================================================
\section{Discussion}
\label{sec:discussion}
% ===================================================================

\subsection{Relationship to optimal experimental design}
\label{sec:design-connection}

The sensor placement problem studied here is a discrete design
problem: the ``design space'' is the set of components
$\{1, \ldots, m\}$, and the experimenter chooses a subset to monitor.
This connects naturally to the optimal experimental design literature
\citep{Atkinson2007, Pukelsheim2006}, where the goal is to choose
design points (here, sensor locations) to optimize a function of the
Fisher information matrix.

The key distinction from classical experimental design is that here the
``design points'' are not freely chosen conditions but structural
features of the system.  The information gain from monitoring a
component depends on the system structure through the critical state
enumeration, creating a non-standard design problem where the design
matrix is determined by the system's reliability logic.

\subsection{Extension to non-exponential lifetimes}
\label{sec:nonexponential}

The monitoring likelihood~\eqref{eq:monitoring-ll} is
distribution-free: it applies to any continuous component lifetime
distributions, not just exponential.  The exponential assumption
enters only in the computation of the density and survival
contributions.  For Weibull or log-normal components, the same
framework applies with modified $f_j$, $F_j$, and $S_j$ functions.

We conjecture that the structural reversal (fastest for parallel,
slowest for voting) persists for other lifetime distributions,
because it depends on the \emph{censoring pattern} imposed by the
system structure, not on the distributional form.  However, the
degree of improvement and the shape of the information gain curve may
differ.

\subsection{Sequential sensor placement}
\label{sec:sequential}

In some applications, sensors can be added sequentially based on
preliminary data.  This suggests an adaptive strategy: estimate
$\blam$ from initial system-level data, identify the information
bottleneck component, and place the first sensor there.  After
collecting monitoring data, re-evaluate and place the next sensor on
the component with the highest remaining marginal gain.  This
sequential approach could be formalized within the Bayesian
experimental design framework \citep{ChalonderVerdinelli1995}.

\subsection{Connection to structural health monitoring}
\label{sec:shm}

The sensor placement problem in reliability systems is structurally
different from the sensor placement problem in structural health
monitoring \citep{OstachowiczSoman2019, Papadimitriou2004}, where
sensors measure continuous signals (vibrations, strains) and
placement is optimized for detecting damage or identifying modal
parameters.  In our setting, sensors measure a binary event (component
failure) and its timing, and placement is optimized for estimating
lifetime distribution parameters.  Despite these differences, both
problems share the feature that the optimal placement depends on the
system's structural properties, and information-theoretic criteria
provide a principled basis for comparison.

\subsection{Limitations}
\label{sec:limitations}

Several limitations should be noted:
\begin{enumerate}
  \item \textbf{Exponential assumption.}  The exponential distribution
    has a constant hazard rate, which may not be realistic for many
    engineering components.  Extension to Weibull and other
    distributions is a natural next step.

  \item \textbf{Small $m$.}  Our simulations use $m = 3$ components.
    The critical state enumeration scales as $O(2^m)$, making
    computation expensive for large~$m$.  Approximations may be needed
    for $m > 10$.

  \item \textbf{Independence assumption.}  We assume component
    lifetimes are independent.  In systems with shared environmental
    stresses or common-cause failures, the optimal sensor placement
    may differ.

  \item \textbf{Limited simulation replications.}  Due to
    computational cost, the $2$-out-of-$3$ experiments use $R = 20$
    replications, which limits the precision of RMSE estimates.

  \item \textbf{$2$-out-of-$3$ all-sensor anomaly.}  The unexpectedly
    high RMSE under full monitoring for the $2$-out-of-$3$ system
    suggests a potential implementation issue that warrants further
    investigation.
\end{enumerate}


% ===================================================================
\section{Conclusion}
\label{sec:conclusion}
% ===================================================================

We have studied the optimal placement of a single sensor for
component lifetime estimation in coherent systems with independent
exponential components.  The central finding is that the optimal
sensor location depends on the system's logical structure through a
structural reversal: in parallel systems, the fastest-failing
component should be instrumented because its lifetime is most often
censored by the system observation, while in voting systems with
$k > 1$, the slowest component should be instrumented because its
status at system failure is most uncertain.

The practical implications are substantial.  For parallel systems, a
single optimally placed sensor reduces the mean RMSE by a factor of
$2.2$ and captures approximately 55--69\% of the full-monitoring
improvement (the lower figure for larger systems, the higher for
three-component systems).  This means that reliability engineers can
achieve most of the benefit of complete component monitoring by
instrumenting a single, judiciously chosen component --- the one with
the highest failure rate in a parallel architecture.

The monitoring likelihood developed in \Cref{sec:likelihood} provides
a general framework for estimating component parameters under partial
monitoring in any coherent system.  Combined with Fisher information
analysis, it enables principled, structure-aware sensor placement
decisions.  Extensions to non-exponential distributions, multiple
sensors, and sequential placement strategies are natural directions
for future work.


% ===================================================================
% Bibliography
% ===================================================================

\bibliographystyle{plainnat}
\bibliography{references}

\end{document}
